Automatic prompt and pipeline optimisation has become routine: a framework searches over instructions, demonstrations and program structure and returns the configuration that scores best on a set of items \citep{khattab2023dspy,opsahlong2024mipro}. The score of that configuration on the items that chose it is the maximum of many noisy scores and is therefore optimistic in the sense of the winner's curse, more so the more candidates were tried and the fewer items were used \citep{cawley2010overfitting,dwork2015preserving}. Good practice already separates tuning, development and test items, and we do not survey how often it is followed. Our subject is the score on whichever items drove the selection: a validation score that is quoted as the result, or a benchmark that was reused for tuning. For that score we ask how much of it survives on fresh items, which factors drive the difference, whether it changes the conclusions drawn from comparisons, and which reporting practice repairs it. The winner's curse itself is not new. What we add is a measurement, against exact truth and with the whole pipeline repeated, of its size for a discrete (correct/incorrect) outcome with correlated candidates and several search procedures; of its effect on three decisions (a tuned-vs-baseline claim, a comparison of systems tuned with unequal budgets, a ranking of tuned systems); and a side-by-side comparison of reporting methods under one reliability criterion.

We call this difference the \emph{optimism gap} $\Delta$ (reported score minus the accuracy of the selected configuration on fresh items) and measure it by repeating the \emph{whole pipeline} thousands of times, so that bias, error and coverage of every reporting method are measured against the truth with Monte-Carlo standard errors. We vary $n$, $K$, the true quality spread, candidate correlation, the search procedure, generation noise and leakage of tuning items into the prompt.

\paragraph{Findings (synthetic benchmark).}
\begin{itemize}\itemsep0pt
\item In the reference cell ($n=200$ tuning items, $K=16$ candidates, true spread $\tau=0.01$) the naive report is $\N{ref_naive_bias}\pm\N{ref_naive_biasse}$~pp above the true accuracy of the selected configuration. The gap exists even when all candidates are exactly equally good ($\tau=0$: $\N{gap_tau00}$~pp at $K=16$) and rises to $\N{gap_K128}$~pp at $K=128$ (reference cell otherwise; Sec.~\ref{sec:synth}).
\item It is driven mainly by $n$ and $K$, is reproduced to within $0.5$~pp by a Gaussian-maximum formula (slightly underpredicted on average), falls with candidate correlation and with the true quality spread, and is raised by leakage and, mildly, by generation noise.
\item It changes conclusions: the naive report licenses ``tuned beats baseline'' in $\N{ref_naive_fp}\%$ of repetitions in which the tuned configuration is truly no better than the baseline (boundary of the null: $5\%$), and systems tuned with unequal budgets are ranked in favour of the larger budget.
\item Holding items out gives an unbiased report (the split meets the reliability criterion in $\N{rel_split_ok}$ of $106$ cells), but at equal total items it has a larger error than an out-of-bag bootstrap correction on all items ($\N{fr_fresh100_200_rmse}$ vs.\ $\N{fr_boot_oob_200_rmse}$~pp RMSE at $200$ items in the reference cell; the bootstrap is better in all $\N{fr_ncells}$ combinations we ran, with coverage $\N{fr_oob_covmin}$--$\N{fr_oob_covmax}\%$). The simultaneous bound under-covers in only $\N{rel_simul_under}$ of $106$ cells but keeps the naive point bias; cross-fitting under-covers. Our implementations of conditional and hybrid winner-aware inference, validated against a closed form and the authors' code, give valid but wide intervals and a noisy point estimate.
\end{itemize}
\paragraph{Findings (real LLM).} Section~\ref{sec:real} repeats the measurement on a small open-weight LLM scoring MMLU-Pro under a pool of prompt configurations. All $K=\N{real_K}$ configurations of a prompt grid are scored (spread $\N{real_sd}$~pp). There the naive report is $\N{real_gap_n200}\pm\N{real_gapse_n200}$~pp above the held-out accuracy at $n=200$ and $\N{real_gap_n50}\pm\N{real_gapse_n50}$~pp at $n=50$ (the same in an i.i.d.\ truth arm), the Gaussian-maximum formula fed with the real covariance reproduces these gaps to about $\N{gm_relmax}\%$, a generator matched to the real matrix reproduces them to within about ten percent, and at equal total items the ordering of the reporting protocols is that of the synthetic study. Successive halving and greedy search shrink it ($\N{real_gap_sh}$, $\N{real_gap_greedy}$~pp) but do not remove it. These are replays over a fixed pool of $\N{real_K}$, not a study of LLM-in-the-loop optimisers or of large $K$.

\paragraph{Scope.} Synthetic statements are about the generator of Sec.~\ref{sec:method}; real-LLM statements are about one small model, one scoring protocol and one item pool. We did not run DSPy or MIPRO: the searches are our re-implementations of their families, and the conditional and hybrid intervals are our implementations of ideas from \citet{andrews2024winners}, checked against a port of the authors' replication code but not identical to it.
